\scanpage{11}

Les structures de similitude strictement fidèles correspondent donc aux \uline{sous-fibrés vectoriels de rang 1} du fibré Vectoriel $\unsure{W|V|\mathbb W}(\underline{\mathrm{Quad}}(E))$, ou encore aux sections du fibré en espaces projectifs $P^\vee(\underline{\mathrm{Quad}}(E))$.

1.4. Soient $\Delta$ un Module inv. sur $X$, et

\[\deleted{\unsure{\Delta}}\qquad\Delta\xrightarrow{\deleted{\unsure{\alpha}}i}\underline{\mathrm{Quad}}(E)\]

un hom (pas néc. injectif), définissant \inserted{donc} une structure de similitude $\Sigma$ sur $E$. Je dis qu'on a une forme quadratique \inserted{$Q_\Sigma$ ou $Q$} canonique à valeurs dans le module inversible $\Delta^\vee$. Pour la définir, je vais d'abord définir (indép. de la donnée de $i$ ou de $\Sigma$) une forme quadratique canonique

\[Q=Q_E:E\longrightarrow\underline{\mathrm{Quad}}(E)^\vee:\]

\deleted{elle associe} sa valeur, pour une section locale $x$ de $E$, est la forme linéaire

\[Q(x)=\{f\longmapsto f(x)\}\quad\inserted{\text{i.e. }f(x)=\langle f,Q(x)\rangle}\]

\inserted{(pour toute section (locale) $f$ de $\underline{\mathrm{Quad}}(E)$, $x$ de $E$,} sur $\underline{\mathrm{Quad}}(E)$. Celle-ci dépend de $x$ de façon quadratique, i.e. la \deleted{fonction} $E\times E\longrightarrow\underline{\mathrm{Quad}}(E)^\vee$,

\[
\unsure{\text{qui transforme}}\quad x,y\ \text{en}\quad Q(x+y)-Q(x)-Q(y)=\{f\longmapsto\underbrace{f(x+y)-f(x)-f(y)}_{\varphi_f(x,y)}\}
\]
% archive sheet number: 10
